\documentclass[11pt,reqno]{amsart}
\usepackage[T1]{fontenc}
\usepackage{lmodern,microtype,amsmath,amssymb,mathtools}
\usepackage[margin=1in]{geometry}
\usepackage[hidelinks]{hyperref}
\usepackage{enumitem,needspace}
\newtheorem{theorem}{Theorem}[section]
\newtheorem{lemma}[theorem]{Lemma}
\newtheorem{proposition}[theorem]{Proposition}
\newtheorem{corollary}[theorem]{Corollary}
\theoremstyle{definition}
\newtheorem{definition}[theorem]{Definition}
\theoremstyle{remark}
\newtheorem{remark}[theorem]{Remark}
\newcommand{\R}{\mathbb R}
\newcommand{\Z}{\mathbb Z}
\newcommand{\E}{\mathbb E}
\newcommand{\1}{\mathbf 1}
\newcommand{\dd}{\,d}
\newcommand{\Var}{\operatorname{Var}}
\newcommand{\Ent}{\operatorname{Ent}}
\newcommand{\Lip}{\operatorname{Lip}}
\newcommand{\Conf}{\operatorname{Conf}}
\newcommand{\cF}{\mathcal F}
\newcommand{\cA}{\mathcal A}
\numberwithin{equation}{section}
\allowdisplaybreaks[2]
\setlength{\emergencystretch}{2em}
\setlist[enumerate]{leftmargin=*,itemsep=4pt}
\title[Particle DLR and stationary uniqueness for 1D jellium]{Particle-only DLR equations and stationary uniqueness\\for one-dimensional Coulomb jellium}
\date{14 September 2026}
\hypersetup{pdftitle={Particle-only DLR equations and stationary uniqueness for one-dimensional Coulomb jellium},pdfsubject={Canonical Coulomb specifications, field reconstruction, and stationary uniqueness}}
\begin{document}
\begin{abstract}
We specify canonical particle-only DLR equations for the genuine
one-dimensional Coulomb potential $g(x)=-|x|$, using symmetric Ces\`aro
averages of particle--background energy differences and an explicit
summability domain. At every inverse temperature $\beta>0$ and background
density $\rho>0$, these equations have exactly one point-process solution
of intensity $\rho$ invariant under all real translations. The principal
step is a more general lifting theorem: finite Gaussian conditional
densities with the Coulomb curvature determine a measurable stationary
electric field through consistency on nested intervals. The exterior
field contains exactly the information in the exterior particles and
the interior count. No discrepancy bound, field moment, or deterministic
asymptotic density is assumed in this theorem. A complete entropy and
Palm proof of stationary field-DLR uniqueness is included. Existence is
deduced from published finite-volume field estimates. The argument does
not require a classification of all nonstationary crystalline phases.
\end{abstract}
\maketitle

\section{Introduction and main conclusions}

The potential considered throughout is
\[
g(x)=-|x|,\qquad -g''=2\delta_0,
\]
with background charge $-\rho\dd x$ and particles of charge $+1$.
This is not the logarithmic interaction, nor the three-dimensional
Coulomb potential restricted to a line. Our pair--background energy
includes a factor $1/2$ in the double integral; its field representation
is $\int E^2$ when $E'=\gamma-\rho\dd x$.

Long-range interactions require a summation convention before an
infinite-volume DLR statement has a definite meaning. Section~\ref{sec:spec}
defines such a convention directly from finite particle--background
energies. It does not put existence of a field lift into the definition.

\begin{theorem}[Stationary particle uniqueness]\label{thm:particle-main}
For every $\beta,\rho>0$, the symmetric Ces\`aro canonical specification
of Definition~\ref{def:particle-dlr} has exactly one probability solution
$P_{\beta,\rho}$ on $\Conf(\R)$ which has intensity $\rho$ and is
invariant under all real translations. Every such solution admits the
deterministic jointly stationary field-DLR lift
\[
Q(\dd\gamma,\dd c)=P_{\beta,\rho}(\dd\gamma)
\delta_{-L_\rho(\gamma)}(\dd c),
\]
where $L_\rho$ is the particle statistic in \eqref{eq:cesaro}.
Its almost-sure density in each direction is $\rho$.
\end{theorem}

There are three separate ingredients. The lifting theorem in
Section~\ref{sec:lift} applies to any particle law with finite,
exterior-measurable Gaussian conditional coefficients. It uses neither
Ces\`aro summability nor a field moment. Sections~\ref{sec:quadratic}--\ref{sec:palm}
prove at most one stationary field-DLR law by a finite-block entropy
comparison. Section~\ref{sec:existence} supplies existence for the
specified normalization using established finite-volume estimates.

The distinction between a prescribed mean intensity and an almost-sure
density is essential: the latter is derived, including for nonergodic
laws. Empty configurations require separate treatment because canonical
resampling can never insert the first particle. They are explicitly
accounted for in Theorem~\ref{thm:lift}.

The classical works of Kunz~\cite{Kunz} and Aizenman--Martin~\cite{AM}
are the relevant historical background. Aizenman--Jansen--Jung~\cite{AJJ}
give an accessible field formulation and the estimates used here.
Our claim is a theorem for the standalone specification stated below;
we make no claim of priority for stationary uniqueness and do not use
a classical classification theorem with an unspecified admissibility class.

\section{Configuration space and field kernels}\label{sec:setup}

Let $\Conf(\R)$ be the space of locally finite simple counting measures,
with its usual Borel structure. Write $N_B(\gamma)=\gamma(B)$ and
$\theta_t\gamma=\gamma-t$. A stationary law of finite intensity
$\lambda$ satisfies $\E N_B=\lambda|B|$ for bounded Borel $B$.
In particular, every deterministic location is unoccupied almost surely.
The background parameter $\rho$ need not initially equal $\lambda$.

For any configuration define the right-continuous anchored primitive
\begin{equation}\label{eq:primitive}
D_\gamma(t)=
\begin{cases}
N_{(0,t]}(\gamma)-\rho t,&t\ge0,\\
-N_{(t,0]}(\gamma)-\rho t,&t<0.
\end{cases}
\end{equation}
Then $D_\gamma(0)=0$, $D_\gamma'=\gamma-\rho\dd x$, and, for $s<t$,
\begin{equation}\label{eq:cocycle}
D_\gamma(t)-D_\gamma(s)=N_{(s,t]}-\rho(t-s).
\end{equation}
All finite-valued right-continuous fields with this derivative have the
form $E_\gamma(x)=C+D_\gamma(x)$.

For $I=(a,b)$ let
\[
\cF_I=\sigma(\gamma_{I^c},N_I).
\]
This is the conditioning sigma-field for a canonical particle kernel.
For $n\ge0$ and $e\in\R$, let $K_{I,n,e}$ be the probability on
ordered points in $I$ with density
\begin{equation}\label{eq:gaussian-kernel}
\frac1{Z_{I,n,e}}\exp\left\{-\beta\rho\sum_{i=1}^n
\left(x_i-a-\frac{e+i-\frac12}{\rho}\right)^2\right\}
\1_{\{a<x_1<\cdots<x_n<b\}}\dd x_1\cdots\dd x_n.
\end{equation}
For $n=0$ this is the point mass on the empty configuration. The
normalizing constant is finite and positive for every finite $e,n$.

\begin{definition}[Canonical field-DLR]\label{def:field-dlr}
A law of $(\gamma,C)$, with $E=C+D_\gamma$, satisfies field-DLR if,
conditionally on the particles and field outside each bounded
deterministic interval $I=(a,b)$, the interior particles have law
$K_{I,N_I,E(a)}$. The exterior field is retained and the interior field is
\begin{equation}\label{eq:dlr}
E(t)=E(a)+\sum_{i=1}^{N_I}\1_{\{x_i\le t\}}-\rho(t-a),
\qquad a<t<b.
\end{equation}
For stationary finite-intensity laws, endpoint particles have probability
zero, and the exterior field determines $N_I=E(b)-E(a)+\rho(b-a)$.
If $0\in I$, $C$ is updated to the value of the resampled field at zero.
\end{definition}

Translations of fields are $E(x)\mapsto E(x+t)$. Joint stationarity is
equivalent to invariance in law under
\[
(\gamma,C)\mapsto(\theta_t\gamma,C+D_\gamma(t)).
\]
Equality of this law for each fixed real $t$ is the stationarity
requirement; a single exceptional null set valid for all $t$ is not
needed.

\section{A particle-only Ces\`aro specification}\label{sec:spec}

\begin{definition}[Summability domain]\label{def:domain}
Let $\cA_\rho$ consist of configurations for which the following limit
exists and is finite:
\begin{equation}\label{eq:cesaro}
L_\rho(\gamma)=\lim_{T\to\infty}\frac1{2T}
\int_{-T}^T D_\gamma(u)\dd u.
\end{equation}
\end{definition}

This is a Borel set: the finite-$T$ expressions are Borel in $\gamma$
and continuous in $T>0$, so rational $T$ suffice to test the limit.
No discrepancy bound or asymptotic density is part of the definition.
The cancellation of the odd function $\rho u$ also shows that the
set and statistic in \eqref{eq:cesaro} are independent of $\rho$;
the background parameter still enters all local energies.

For $I=(a,b)$, an exterior configuration $\zeta$, and an interior
configuration $\xi$, write
\[
\mu_R^\xi=(\xi+\zeta-\rho\dd x)|_{(-R,R]},\qquad
U_R(\xi)=\frac12\iint g(x-y)\dd\mu_R^\xi(x)\dd\mu_R^\xi(y).
\]
Endpoint conventions for $R$ do not change any of the averages below.
There is no point self-energy divergence since $g(0)=0$.
For two interiors $\xi,\eta$ of equal cardinality define
\begin{equation}\label{eq:relative-energy}
\Delta H_I(\xi,\eta;\zeta)
=\lim_{T\to\infty}\frac1T\int_{R_0}^T
\bigl(U_R(\xi)-U_R(\eta)\bigr)\dd R,
\end{equation}
where $(-R_0,R_0)$ contains $\overline I$.

\begin{proposition}[Finite local energies and canonical closure]\label{prop:energy}
Suppose $\eta+\zeta\in\cA_\rho$ and $|\xi|=|\eta|$. Then
$\xi+\zeta\in\cA_\rho$, \eqref{eq:relative-energy} exists, and,
on writing
\begin{equation}\label{eq:normalized-field}
\mathcal E_\gamma(x)=D_\gamma(x)-L_\rho(\gamma),
\end{equation}
one has
\begin{align}
\mathcal E_{\xi+\zeta}&=\mathcal E_{\eta+\zeta}
&&\text{outside }I,\label{eq:retained-field}\\
\Delta H_I(\xi,\eta;\zeta)
&=\int_a^b\left(\mathcal E_{\xi+\zeta}^2
-\mathcal E_{\eta+\zeta}^2\right)\dd t.\label{eq:energy-field}
\end{align}
The relative energy is independent of $R_0$ and satisfies the cocycle
identity in its two interior arguments.
\end{proposition}

\begin{proof}
Let $\delta=\xi-\eta$ and let $d(t)=\delta(( -\infty,t])$.
Since $\delta(\R)=0$, $d$ vanishes outside $I$. The anchored
primitives satisfy
\[
D_{\xi+\zeta}(t)-D_{\eta+\zeta}(t)=d(t)-d(0).
\]
Taking symmetric means gives
$L_\rho(\xi+\zeta)-L_\rho(\eta+\zeta)=-d(0)$.
This proves canonical closure and
$\mathcal E_{\xi+\zeta}-\mathcal E_{\eta+\zeta}=d$.

For $|t|<R$, the right-continuous field $-\frac12(g*\mu_R^\eta)'$
is
\begin{equation}\label{eq:cutoff-field}
E_R^\eta(t)=D_{\eta+\zeta}(t)
-\frac{D_{\eta+\zeta}(R)+D_{\eta+\zeta}(-R)}2.
\end{equation}
The integration-by-parts identity for the neutral compactly supported
measure $\delta=d'$ gives
\[
U_R(\xi)-U_R(\eta)=\int_I\left(2E_R^\eta(t)d(t)+d(t)^2\right)\dd t.
\]
No neutrality assumption on $\mu_R^\eta$ is necessary: $d$ is compactly
supported, so no boundary term remains. The Ces\`aro mean of the last
term in \eqref{eq:cutoff-field} is $L_\rho(\eta+\zeta)$, which proves
\eqref{eq:energy-field}. The remaining statements follow directly.
\end{proof}

\begin{definition}[Symmetric Ces\`aro canonical particle-DLR]\label{def:particle-dlr}
For $\gamma\in\cA_\rho$, retain $\zeta=\gamma_{I^c}$ and $n=N_I$.
The conditional density of $n$ ordered interior points is proportional
to
\begin{equation}\label{eq:particle-dlr}
\exp[-\beta\Delta H_I(\xi,\eta;\zeta)]
\1_{\{a<x_1<\cdots<x_n<b\}}\dd x_1\cdots\dd x_n,
\end{equation}
where $\eta$ is any $n$-point completion in $I$. A probability $P$
satisfies these DLR equations if $P(\cA_\rho)=1$ and
\eqref{eq:particle-dlr} is its conditional law given $\cF_I$ for
every bounded deterministic interval $I$.
\end{definition}

The definition is standalone and uses only particle--background
energies. Formula \eqref{eq:normalized-field} is a derived statistic,
not an assumption about a jointly stationary extension.

For explicit measurability, choose the reference completion
$\eta=\sum_{i=1}^n\delta_{a+i(b-a)/(n+1)}$. Canonical closure shows
that admissibility can be tested on this completion. The parameter
\begin{equation}\label{eq:particle-coefficient}
e_I(\zeta,n)=\mathcal E_{\eta+\zeta}(a)
\end{equation}
is therefore a Borel function of the retained data. It agrees with
$\mathcal E_\gamma(a)$ for any admissible completion.
Moving the $i$th ordered jump gives
\[
\partial_{x_i}\int_I\mathcal E^2
=2\rho(x_i-a)-2e_I-2i+1.
\]
Thus \eqref{eq:particle-dlr} is exactly $K_{I,n,e_I}$. In particular,
its normalization is finite and positive. The kernels retain
admissibility and are consistent under nesting: their exterior field
is fixed by \eqref{eq:retained-field}, and the integral energy splits
over subintervals. These statements also follow by conditioning the
explicit ordered Gaussian density.

\begin{remark}[Scope of the normalization]\label{rem:scope}
On arbitrary configurations $\cA_\rho$ need not be invariant under
translation, and the normalized primitive need not be covariant.
We assert neither property. They hold almost surely in the sense needed
for stationary candidate laws by Theorem~\ref{thm:lift} below.
The zero symmetric mean of $\mathcal E$ is a consequence of the declared
summation convention, not a field-moment assumption or a consequence
claimed from mean intensity alone.

The summability issue is intrinsic to \eqref{eq:relative-energy}:
moving one particle so that $\int_I d\ne0$ makes the averaged energy
difference a constant minus a nonzero multiple of the symmetric average
in \eqref{eq:cesaro}. Thus finite local-move energy differences for this
convention require that limit. Alternative renormalizations must be
specified separately.
\end{remark}

\section{Lifting finite Gaussian particle kernels}\label{sec:lift}

The following theorem is independent of Section~\ref{sec:spec}.
It is the link between particle-only conditional laws and a stationary
field formulation.

\begin{theorem}[Lifting theorem]\label{thm:lift}
Let $P$ be a stationary law on $\Conf(\R)$ of finite positive intensity
$\lambda$. Suppose that for every bounded deterministic interval
$I=(a,b)$ there is a finite real $\cF_I$-measurable variable $e_I$ such
that, on $\{N_I=n>0\}$, the conditional interior law is
$K_{I,n,e_I}$. No condition on $e_I$ on $\{N_I=0\}$ is needed.
Then, conditional on nonemptiness, $P$ admits a measurable deterministic
jointly stationary field-DLR lift $E=C(\gamma)+D_\gamma$.
The conditional law on nonempty configurations has intensity and
almost-sure forward and backward density $\rho$. Consequently
\begin{equation}\label{eq:empty-mass}
P(\varnothing)=1-\lambda/\rho,
\end{equation}
so necessarily $\lambda\le\rho$. In particular, if $\lambda=\rho$,
the whole law lifts. For every fixed $t\in\R$,
\begin{equation}\label{eq:lift-covariance}
C(\theta_t\gamma)=C(\gamma)+D_\gamma(t)\qquad P\text{-a.s.}
\end{equation}
on the nonempty component.
\end{theorem}

\begin{proof}
The event $\{\gamma\ne\varnothing\}$ is measurable from $\cF_I$
for every $I$: either the retained exterior is nonempty or $N_I>0$.
Conditioning on this event preserves all the conditional laws and
stationarity. Its probability is positive and its conditional intensity
is finite. Work on that component until the final paragraph.

\emph{Compatibility and identifiability.}
Let $I=(a,b)\subset J=(c,d)$. Then $\cF_J\subset\cF_I$.
Further conditioning the $J$-density on the exterior of $I$ fixes the
number $k=N_{(c,a]}$ of particles to the left of $I$. The remaining
ordered density in $I$ is Gaussian with coefficient
$e_J+k-\rho(a-c)$. This follows by replacing the rank $i$ in the
$I$-block with $k+i$ in the $J$-block. Disintegration is legitimate
because the $J$-density is strictly positive on its ordered chamber.

Two coefficients $e,e'$ in $K_{I,n,e}$, $n\ge1$, have likelihood ratio
\[
\text{constant}\times\exp\left(2\beta(e-e')\sum_{i=1}^n x_i\right).
\]
This ratio can be constant almost everywhere on the chamber only if
$e=e'$. Uniqueness of conditional laws therefore gives
\begin{equation}\label{eq:nested-coefficient}
e_I=e_J+N_{(c,a]}-\rho(a-c)
\quad\text{on }\{N_I>0\},\quad P\text{-a.s.}
\end{equation}

\emph{Measurable patching.}
Take all rational bounded intervals, together with countably many
rational enclosing intervals as needed. Their identities
\eqref{eq:nested-coefficient} hold simultaneously outside one null set.
For occupied intervals, \eqref{eq:cocycle} implies that
\begin{equation}\label{eq:patch-C}
e_I-D_\gamma(a)
\end{equation}
has the same value for any two choices of $I$: compare both to a common
occupied enclosing interval. Select the first occupied member of the
sequence $(-m,m)$, $m\ge1$, and use \eqref{eq:patch-C} to define $C$.
This is finite and measurable. Define $E=C+D_\gamma$.
For any additional fixed interval, the same identity follows by
comparison to the exhausting rational sequence.

\emph{The conditioning sigma-fields.}
Fix $I=(a,b)$, whether or not it is occupied. Take a deterministic
increasing sequence of intervals $K_m=(u_m,v_m)\supset I$ exhausting
$\R$ and choose its first occupied member $K$. This choice is
$\cF_I$-measurable, since $N_{K_m}=N_I+N_{K_m\setminus I}$.
Both $e_K$ and the count $N_{(u_K,a]}$ are $\cF_I$-measurable.
Patching gives
\begin{equation}\label{eq:outside-measurability}
E(a)=e_K+N_{(u_K,a]}-\rho(a-u_K).
\end{equation}
Every exterior field value follows from this boundary value and the
exterior counts, adding the known total $N_I$ when crossing $I$.
Hence the exterior field is $\cF_I$-measurable. Conversely it recovers
$N_I$ from $E(b)-E(a)+\rho(b-a)$. It follows that
\begin{equation}\label{eq:sigma-equality}
\sigma(\gamma_{I^c},E|_{I^c})=\cF_I\quad\text{modulo null sets}.
\end{equation}
Right continuity permits use of countably many field evaluations here;
the left boundary value is also recovered by a limit from the exterior.
Deterministic endpoints are unoccupied almost surely.

On $N_I>0$, \eqref{eq:nested-coefficient} identifies $E(a)=e_I$.
On $N_I=0$ the conditional configuration is empty. Thus
\eqref{eq:sigma-equality} proves field-DLR. Formula
\eqref{eq:outside-measurability} also shows that exterior field values
are retained under conditional resampling. When $0\in I$, the new
value of $C$ is the value at zero of the bridge \eqref{eq:dlr}, so it
is updated rather than retained.

\emph{Stationarity of the lift.}
For each fixed real $t$, stationarity of $P$ makes the translated
$I+t$ conditional law another version of the $I$ conditional law.
Identifiability yields
\[
e_I(\theta_t\gamma)=e_{I+t}(\gamma)
\quad\text{on }\{N_{I+t}>0\},\quad P\text{-a.s.}
\]
Compare the countable rational family, its $t$-translate, and a
countable family of enclosing intervals on one full-measure set.
Since $D_{\theta_t\gamma}(a)=D_\gamma(a+t)-D_\gamma(t)$,
\eqref{eq:patch-C} gives \eqref{eq:lift-covariance}. This proves joint
stationarity, without asserting an uncountable intersection of null-set
complements.

\emph{Density and the empty component.}
Stationarity of the finite-valued field gives
\[
\frac{N_{(0,t]}}t-\rho=\frac{E(t)-E(0)}t\longrightarrow0
\quad\text{in probability}.
\]
Indeed, each of $E(t)$ and $E(0)$ has the same fixed tight distribution.
The point-process ergodic theorem at finite intensity gives almost-sure
and $L^1$ convergence of $N_{(0,t]}/t$ to the invariant-component
intensity. It must therefore equal $\rho$. The backward argument is
identical. In particular the nonempty component has mean intensity
$\rho$. If $p=P(\varnothing)$ for the original law, its intensity is
$\lambda=(1-p)\rho$, proving \eqref{eq:empty-mass}.
\end{proof}

\begin{corollary}[Application to the particle specification]\label{cor:cesaro-lift}
Every stationary solution of Definition~\ref{def:particle-dlr} of
intensity $\rho$ has the lift $C=-L_\rho(\gamma)$ in
Theorem~\ref{thm:particle-main}.
\end{corollary}
\begin{proof}
The coefficient \eqref{eq:particle-coefficient} is finite and
$\cF_I$-measurable. Apply Theorem~\ref{thm:lift}; on any occupied
interval, \eqref{eq:patch-C} equals $-L_\rho(\gamma)$.
\end{proof}

\begin{remark}[Which other specifications are covered?]
The theorem concerns conditional densities, not names of normalization
schemes. For $g=-|x|$, an exterior particle potential is affine on $I$;
the local background has second derivative $2\rho$, and the ordered
pair energy is linear. Consequently every finite real exterior-potential
prescription of this form yields \eqref{eq:gaussian-kernel} and is
covered once its measurability and admissibility have been checked.
A kernel obtained by mixing distinct unobserved field coefficients
need not have this form. Neither such mixtures nor kernels singular
with respect to ordered Lebesgue measure are included by assertion.
\end{remark}

\section{Stationary field uniqueness and quadratic coordinates}\label{sec:quadratic}

\begin{theorem}[Stationary field-DLR uniqueness]\label{thm:field-unique}
Fix $\beta,\rho>0$. There is at most one jointly stationary law of
$(\gamma,C)$ of particle intensity $\rho$ with finite-valued
$E=C+D_\gamma$ satisfying Definition~\ref{def:field-dlr}.
No field-moment assumption is required.
\end{theorem}

We give the proof in full in this and the next three sections.



For $a<x_1<\cdots<x_n<b$, write
\[
 H_{a,b,e}(x)=\int_a^b
 \left(e+\sum_{i=1}^n\1_{\{x_i\le t\}}-\rho(t-a)\right)^2\dd t.
\]
Moving the $i$th jump gives
\begin{align*}
 \partial_{x_i}H_{a,b,e}
 &=\big(e+i-1-\rho(x_i-a)\big)^2
   -\big(e+i-\rho(x_i-a)\big)^2\\
 &=2\rho(x_i-a)-2e-2i+1.
\end{align*}
Consequently, on the ordered chamber,
\begin{equation}\label{eq:quadratic}
 H_{a,b,e}(x)=\rho\sum_{i=1}^n
 \left(x_i-a-\frac{e+i-\frac12}{\rho}\right)^2+C(a,b,e,n),
 \qquad H_{a,b,e}''=2\rho I_n.
\end{equation}
No neutrality relation $n=\rho(b-a)$ is required here. The background
provides the entire Hessian; the ordered pair interaction is linear.

Set
\begin{equation}\label{eq:y}
 y_i=\rho(x_i-a)-e+\tfrac12,\qquad \alpha=\beta/\rho.
\end{equation}
Let $\nu_n=\nu_{n,\alpha}$ be the probability measure on the unbounded
ordered chamber $\mathcal O_n=\{y_1<\cdots<y_n\}$ with density
\begin{equation}\label{eq:nu}
 \nu_n(\dd y)=Z_n^{-1}
 \exp\left[-\alpha\sum_{i=1}^n(y_i-i)^2\right]\dd y.
\end{equation}
Then the conditional law of $y$ is precisely
\begin{equation}\label{eq:walls}
 \mu_{n;A,B}=\nu_n(\,\cdot\mid W_{A,B}),\qquad
 W_{A,B}=\{A<y_1,\ y_n<B\},
\end{equation}
where
\begin{equation}\label{eq:AB}
 A=\tfrac12-e,\qquad B-n=\tfrac12-f,
 \qquad f=e+n-\rho(b-a).
\end{equation}
For particle endpoints, the same formulas hold with $e=E(a+)$ and
$f=E(b-)$, since the endpoints themselves are retained.

There are two useful cancellation identities:
\begin{equation}\label{eq:marks}
 y_{i+1}-y_i=\rho(x_{i+1}-x_i),\qquad
 z_i:=y_i-i=\tfrac12-E(x_i+).
\end{equation}
Thus the unknown common shift of the centers disappears from physical
gaps, while the displacements $z_i$ encode field marks exactly.

\section{The entropy cost of the walls}\label{sec:wall}

\begin{lemma}[A uniform sublinear wall cost]\label{lem:wall}
For every $\alpha>0$ and $M<\infty$,
\begin{equation}\label{eq:wall-small}
 \sup_{|A|\le M,\ |B-n|\le M}
 \frac{-\log\nu_n(W_{A,B})}{n}\longrightarrow0.
\end{equation}
More explicitly, for all sufficiently large $n$ depending on $\alpha,M$,
\begin{equation}\label{eq:wall-rate}
 -\log\nu_n(W_{A,B})
 \le C_\alpha\big(M+2+\sqrt{\log(4n)/\alpha}\big)n^{2/3}.
\end{equation}
\end{lemma}

\begin{proof}
We first bound the two extreme coordinates. Let $\sigma_c$ be the law
of the ordered values of $n$ independent Gaussians of mean $c$ and
variance $(2\alpha)^{-1}$. This law is positively associated: sorting is
coordinatewise increasing, and independent variables are associated.
On $\mathcal O_n$,
\[
 \frac{\dd\nu_n}{\dd\sigma_c}(y)
 =\text{constant}\times
 \exp\left(2\alpha\sum_{i=1}^n(i-c)y_i\right).
\]
For $c=1$ the ratio is increasing in every coordinate, so association
implies $\sigma_1\preceq\nu_n$ in stochastic order. For $c=n$ the ratio
is decreasing, and $\nu_n\preceq\sigma_n$. The covariance inequalities
extend to these integrable exponential ratios by truncation. Gaussian
tail bounds and a union bound therefore give, for $R>0$,
\begin{equation}\label{eq:extreme}
 \nu_n(y_1<-R)+\nu_n(y_n>n+R)
 \le 2n e^{-\alpha R^2}.
\end{equation}
In particular, with $R=\sqrt{\log(4n)/\alpha}$, the event
\[
 C_n=\{y_1\ge -R,\ y_n\le n+R\}
\]
has probability at least $1/2$.

Put $r=R+M+2$, $m=\lceil n^{1/3}\rceil$, and
$\varepsilon=rn^{-1/3}$. For sufficiently large $n$, we have
$\varepsilon<1/2$ and $n>2m+1$. Define
\[
 d_i=r\left(1-\frac{i-1}{m}\right)_+
     -r\left(1-\frac{n-i}{m}\right)_+,
 \qquad
 T_i(y)=(1-\varepsilon)y_i+\varepsilon i+d_i.
\]
The two supports in $d_i$ are disjoint, and
$d_{i+1}-d_i\ge-r/m\ge-\varepsilon$. Hence
\[
 T_{i+1}(y)-T_i(y)
 \ge(1-\varepsilon)(y_{i+1}-y_i)>0.
\]
For $y\in C_n$,
\[
 T_1(y)\ge M+2+\varepsilon(R+1)>A,
 \qquad
 T_n(y)\le n-M-2-\varepsilon R<B.
\]
Thus $T(C_n)\subset W_{A,B}$. This affine map is injective and has
Jacobian $(1-\varepsilon)^n$.

Writing $q_i=y_i-i$, convexity of the square gives
\[
 (T_i(y)-i)^2=((1-\varepsilon)q_i+d_i)^2
 \le(1-\varepsilon)q_i^2+\frac{d_i^2}{\varepsilon}
 \le q_i^2+\frac{d_i^2}{\varepsilon}.
\]
Also $\sum_i d_i^2\le2(m+1)r^2$. Changing variables on $T(C_n)$ yields
\[
 \nu_n(W_{A,B})
 \ge\frac12(1-\varepsilon)^n
 \exp\left[-\frac{2\alpha(m+1)r^2}{\varepsilon}\right].
\]
Since $-\log(1-\varepsilon)\le2\varepsilon$,
\[
 -\log\nu_n(W_{A,B})
 \le\log2+2n\varepsilon+
       \frac{2\alpha(m+1)r^2}{\varepsilon}
 \le C_\alpha r n^{2/3}.
\]
This proves both claims.
\end{proof}

\section{Gaussian entropy comparison for block averages}\label{sec:comparison}

Write $S(\mu\mid\nu)=\int\log(\dd\mu/\dd\nu)\dd\mu$.
The Hessian of the exponent of \eqref{eq:nu} is $2\alpha I_n$.
The logarithmic Sobolev inequality on a convex domain and its Herbst
consequence imply, for an $L$-Lipschitz $F$,
\[
 \log\nu_n\!\left(e^{t(F-\nu_nF)}\right)
 \le\frac{t^2L^2}{4\alpha},\qquad t\in\R.
\]
For the concentration and logarithmic Sobolev inequalities, see
\cite{BGL}. They apply here by approximating the indicator of the
closed ordered chamber by smooth convex penalties, retaining Hessian
at least $2\alpha I_n$, and passing to the limit. The boundary has
zero Lebesgue measure; the limiting density need not vanish there. Applying the entropy variational inequality with either sign
and optimizing $t$ gives
\begin{equation}\label{eq:entropy}
 |\mu F-\nu_nF|^2\le\frac{L^2}{\alpha}S(\mu\mid\nu_n).
\end{equation}

For a fixed integer $r\ge1$ and bounded Lipschitz $G:\R^r\to\R$, define
the displacement-mark average
\begin{equation}\label{eq:mark-stat}
 F_n(y)=\frac1{n-r+1}\sum_{i=1}^{n-r+1}
 G(y_i-i,\ldots,y_{i+r-1}-(i+r-1)).
\end{equation}
Bounded overlap of windows gives the explicit bound
$\Lip(F_n)^2\le r\Lip(G)^2/(n-r+1)$, hence
$\Lip(F_n)^2\le C_{G,r}/n$ for $n\ge2r$. The same estimate holds for
averages of bounded Lipschitz functions of any fixed number of
consecutive gaps $y_{i+1}-y_i$.

For the conditioned law in \eqref{eq:walls}, exactly
\begin{equation}\label{eq:entropy-exact}
 S(\mu_{n;A,B}\mid\nu_n)=-\log\nu_n(W_{A,B}).
\end{equation}
Combining \eqref{eq:entropy}, \eqref{eq:entropy-exact}, and
Lemma~\ref{lem:wall} proves
\begin{equation}\label{eq:uniform-comparison}
 \sup_{|A|\le M,\ |B-n|\le M}
 |\mu_{n;A,B}F_n-\nu_nF_n|\longrightarrow0.
\end{equation}
In particular, its square is bounded by
$C_{\alpha,G,r}(M+2+\sqrt{\log(4n)/\alpha})n^{-1/3}$ for large $n$.

\begin{remark}[Why use entropy directly?]
The wall indicator makes a direct relative-Fisher-information comparison
with the unbounded reference generally infinite. Equation
\eqref{eq:entropy-exact} is finite and is the appropriate replacement.

\end{remark}

\section{Stationary random blocks and Palm identification}\label{sec:palm}

Let $P$ be a law in Theorem~\ref{thm:field-unique}, and write $P^0$ for its
marked Palm law. First, the density is $\rho$ almost surely in both
directions. Indeed, stationarity of the finite-valued field gives
\[
 \frac{N_{(0,t]}}t-\rho=\frac{E(t)-E(0)}t\longrightarrow0
 \quad\hbox{in probability}.
\]
The stationary ergodic theorem for the finite-intensity point process
gives an almost-sure count-density limit, which must therefore be
$\rho$, even if $P$ is not ergodic. The backward argument is identical.
This translation-invariant property passes to Palm measure.

Under $P^0$ enumerate the points by
$\cdots<p_{-1}<p_0=0<p_1<\cdots$, put $g_i=p_{i+1}-p_i$, and let
$T_*$ be the shift to the successor, translating the field as well.
Marked Palm measure is $T_*$-invariant and $\E_{P^0}g_{-1}=1/\rho$.
For a spatial sample, let $a$ be its first positive particle.
The law of the marked environment viewed from $a$ is
\begin{equation}\label{eq:bias}
 \widetilde P^0(\dd\omega)=\rho g_{-1}(\omega)P^0(\dd\omega).
\end{equation}
Indeed, Palm inversion over the preceding gap gives
$\E_P h(\theta_a\omega)=\rho\E_{P^0}[g_{-1}h]$ for bounded
measurable $h$. This proves the displayed length-bias formula; see
\cite[Chapter~9]{LastPenrose} for Palm inversion.

Let $b$ be the $(n+2)$nd positive particle, so that $(a,b)$ contains
exactly $n$ points $x_1<\cdots<x_n$. Retain $a,b$, every other exterior
particle, and the field outside $(a,b)$.

\begin{lemma}[Random-block disintegration and tight walls]\label{lem:block}
Conditionally on this retained information, the transformed interior
coordinates \eqref{eq:y}, with $e=E(a+)$, have law $\mu_{n;A,B}$ in
\eqref{eq:walls}. The family $(A,B-n)$ is tight as $n\to\infty$.
\end{lemma}

\begin{proof}
For disintegration, first restrict to $b<R$ for a deterministic integer
$R$. Apply field-DLR in $(0,R)$ and condition further on every ordered
coordinate except the $n$ middle particles, retaining the resulting
exterior field. Moving these middle particles within $(a,b)$ preserves
the selected ranks. The part of $\int_0^R E^2$ outside $(a,b)$ is fixed.
Thus the remaining density is \eqref{eq:dlr} on $(a,b)$. The events
$\{b<R\}$ are retained-information measurable and exhaust probability
one. Equations~\eqref{eq:quadratic}--\eqref{eq:AB} identify the law.

To prove tightness, let $g_*$ be the entire gap containing the spatial
origin. On $\{g_*\le L\}$, the density in \eqref{eq:bias} is bounded
by $\rho L$. Invariance under $T_*$ gives, uniformly in $n$,
\begin{align*}
 &P\big(|E(a+)|>R\ \hbox{or}\ |E(b-)|>R\big)\\
 &\quad\le P(g_*>L)+\rho L\left\{
 P^0(|E(0+)|>R)+P^0(|E(0-)|>R)\right\}.
\end{align*}
The Palm field values are finite almost surely. First choose $L$ large,
then $R$ large. Formula~\eqref{eq:AB} proves the claim. No Palm moment
of the field is being assumed.
\end{proof}

\begin{proof}[Proof of Theorem~\ref{thm:field-unique}]
Fix $r$ and a bounded Lipschitz $G$ as in \eqref{eq:mark-stat}.
By \eqref{eq:marks}, the physical block statistic is
\[
 \mathfrak F_n(\gamma,E)=\frac1{n-r+1}\sum_{i=1}^{n-r+1}
 G\big(\tfrac12-E(x_i+),\ldots,
       \tfrac12-E(x_{i+r-1}+ )\big).
\]
Its conditional expectation is $\mu_{n;A,B}F_n$.
Lemma~\ref{lem:block}, \eqref{eq:uniform-comparison}, and boundedness
give
\begin{equation}\label{eq:limit-reference}
 \E_P\mathfrak F_n-\nu_nF_n\longrightarrow0.
\end{equation}
Explicitly, restrict first to $|A|,|B-n|\le M$, take $n\to\infty$,
and then let $M\to\infty$ using uniform tightness.

We identify the limit on the left via marked Palm measure. Let
$\mathcal I$ be the $T_*$-invariant sigma-algebra. Birkhoff's theorem
for the bounded mark observable $G$ yields the conditional limit
$\E_{P^0}[G\mid\mathcal I]$ for the successive block averages.
The derived spatial density implies $p_m/m\to1/\rho$ under Palm
measure. Birkhoff applied to the integrable gap therefore gives
\[
 \E_{P^0}[g_{-1}\mid\mathcal I]=1/\rho.
\]
Use \eqref{eq:bias}, dominated convergence with weight $\rho g_{-1}$,
and conditioning on $\mathcal I$. It follows that
\begin{equation}\label{eq:palm-limit}
 \E_P\mathfrak F_n\longrightarrow
 \E_{P^0}G\big(\tfrac12-E(p_0+),\ldots,
              \tfrac12-E(p_{r-1}+ )\big).
\end{equation}
The finitely shifted starting index does not change this limit.

For any second candidate $Q$, the same deterministic $\nu_nF_n$ occurs
in \eqref{eq:limit-reference}. Equations
\eqref{eq:limit-reference}--\eqref{eq:palm-limit} therefore identify
all finite consecutive distributions of
$z_i=\tfrac12-E(p_i+)$ under $P^0$ and $Q^0$. Successor invariance
identifies every finite window with arbitrary integer indices.

The entire marked Palm configuration is determined by this sequence:
\begin{equation}\label{eq:reconstruct}
 \rho g_i=1+z_{i+1}-z_i,\qquad p_0=0,
 \qquad E(p_i+)=\tfrac12-z_i,
\end{equation}
and the field has slope $-\rho$ between particles. Thus the marked
Palm laws coincide. Marked Palm inversion, for bounded measurable $h$,
\[
 \E_Ph=\rho\E_{P^0}\int_0^{g_0}
 h\big(\theta_t(\gamma,E)\big)\dd t,
\]
now proves $P=Q$.
\end{proof}




\section{Existence for the specified normalization}\label{sec:existence}

Only existence uses estimates from the classical finite-volume
construction.  The preceding lifting and uniqueness arguments apply to
every candidate in their stated classes.

For $N\in\mathbb N$, let $P_N$ be the law of $2N$ ordered particles in
$B_N=(-N/\rho,N/\rho)$, with density proportional to
\[
 \exp\left[-\beta\int_{B_N} E_N(x)^2\,dx\right],\qquad
 E_N(x)=N_{(-N/\rho,x]}-\rho(x+N/\rho),\quad x\in B_N.
\]
The field has value zero at both endpoints. Extend it to the whole line
with derivative $\gamma_N-\rho\,dx$, where $\gamma_N$ has no particles
outside $B_N$. This extension is used only to put all laws on one state
space; finite-volume energies are integrated over $B_N$.

\begin{lemma}[Imported finite-volume estimates]\label{lem:fv-estimates}
For fixed $\beta,\rho>0$, there are positive constants $A,B$ such that,
uniformly in $N$ and $x\in B_N$,
\begin{equation}\label{eq:fv-tail}
 P_N(|E_N(x)|>u)\le A e^{-Bu^3},\qquad u\ge0.
\end{equation}
For every $\delta>0$, there are $A_\delta,B_\delta>0$ such that, uniformly
over intervals $J\subset B_N$,
\begin{equation}\label{eq:fv-average}
 P_N\left(\left|\frac1{|J|}\int_J E_N(x)\,dx\right|>\delta\right)
 \le A_\delta e^{-B_\delta |J|}.
\end{equation}
\end{lemma}

These are the strictly one-dimensional estimates of
\cite[Theorem~6.2 and Proposition~8.1]{AJJ}; the latter is the
one-dimensional proof of their volume-average bound, Theorem~3.4.
Their unit-density convention is $U=\frac12\int K^2$.
The conversion is $y=\rho x$, $K(y)=-E_N(y/\rho)$, and
$\beta_{\mathrm{AJJ}}=2\beta/\rho$. No numerical value of an exponential
constant is needed here. The displayed bounds at arbitrary spatial
endpoints follow from the bounds on the density grid: on each grid cell,
the field is bounded in absolute value by one plus the larger absolute
endpoint value. Rounding the endpoints of $J$ changes its integral only
by two such cell integrals; \eqref{eq:fv-tail} controls those errors.
Thus the constants are uniform in the box and the interval.

\begin{proposition}\label{prop:existence}
There exists a jointly stationary field-DLR law of intensity $\rho$.
It has $\E|E(0)|<\infty$ and almost surely
\begin{equation}\label{eq:constructed-zero-mean}
 \lim_{R\to\infty}\frac1{2R}\int_{-R}^{R}E(x)\,dx=0.
\end{equation}
Its particle marginal satisfies Definition~\ref{def:particle-dlr}.
The field has finite moments of every finite order.
\end{proposition}

\begin{proof}
Independently choose $U_N$ uniformly in
$[-N/(2\rho),N/(2\rho)]$ and let $\overline Q_N$ be the law of
$(\theta_{U_N}\gamma_N,E_N(\,\cdot+U_N))$.
Use vague convergence for locally finite integer-valued measures
(temporarily allowing multiple points), and $L^1_{\mathrm{loc}}$
convergence for fields. Translations are continuous in this topology.
Changing $U_N$ to $U_N+t$ changes its law in total variation by
$O(\rho|t|/N)$. Consequently every subsequential limit is stationary.

Here is the required compactness argument. For every fixed compact
interval $K$, all its shifted points lie inside $B_N$ for large $N$.
Equation~\eqref{eq:fv-tail} bounds the expectation of the field's
$L^1(K)$ norm uniformly. For fixed $a<b$ in that region,
\[
 N_{(a,b]}=E_N(b)-E_N(a)+\rho(b-a).
\]
The same estimate therefore bounds the expected number of particles
on $K$. Since the distributional derivative is
$\gamma_N-\rho\,dx$, this also bounds the expected total variation
on $K$. Compactness of bounded sets of $BV(K)$ in $L^1(K)$, followed
by a diagonal argument and the corresponding vague compactness for
counting measures, proves tightness. Choose a limit $Q$.
The distributional identity passes to the limit, so
$E'=\gamma-\rho\,dx$. The field admits a unique right-continuous
representative and is finite-valued. The limit counting measure remains
integer-valued, with possible multiplicities at this stage.

Stationarity and the count bounds imply that its mean counting measure
is a finite constant times Lebesgue measure. In particular, no
deterministic point is an atom almost surely. The local $L^1$ bound
passes to the limit; stationarity then gives $\E_Q|E(0)|<\infty$.
Taking expectations in the charge identity proves
$\E_Q N_{(a,b]}=\rho(b-a)$.

We next pass the conditional equations, spelling out the continuity
point that prevents a hidden simplicity assumption. For a fixed
deterministic $I=(a,b)$ let $\Gamma_I$ replace its particles by the
Gaussian density in Definition~\ref{def:field-dlr} and retain its
exterior field. Define this kernel also for integer-valued incoming
measures with multiple points. If the incoming measure has no atoms at
$a,b$, its interior count is locally constant under vague convergence,
and its boundary field value converges under the joint topology above.
The latter assertion follows from $E'=\gamma-\rho\,dx$: convergence
of monotone primitives holds at their continuity points.
For each fixed count, the ordered Gaussian density depends
continuously on the finite boundary value, by dominated convergence on
the bounded chamber. Hence $\Gamma_I f$ is continuous there for every
bounded continuous $f$ on the joint state space. Resampled fields also
converge in $L^1_{\mathrm{loc}}$.

For large $N$, every shifted box contains $I$, and the finite-volume
conditional equation holds for each origin in the averaging range.
Thus $\overline Q_N\Gamma_I=\overline Q_N$. The almost-sure continuity
just established gives $Q\Gamma_I=Q$. As $\Gamma_I$ retains the
exterior information, this identity is equivalent to the conditional
field-DLR equation. Moreover, $\Gamma_I$ produces distinct interior
points with probability one. Applying this to a countable covering of
the line proves that $\gamma$ is simple almost surely.
The boundary-evaluation continuity also transfers \eqref{eq:fv-tail}
to the limiting field at each deterministic point (with an arbitrarily
smaller threshold before passage to the limit). Thus the limiting field
has finite moments of every finite order.

For each fixed $R$, \eqref{eq:fv-average} passes to the averaged laws
and their limit, because integration on $[-R,R]$ is continuous in
$L^1_{\mathrm{loc}}$. It bounds the probability that the absolute
average exceeds $\delta$ by
$A_\delta e^{-2B_\delta R}$. Borel--Cantelli gives zero limiting average
along integer $R$. The stationary ergodic theorem, applicable because
$\E_Q|E(0)|<\infty$, gives an almost-sure limit for the full real
parameter; that limit must be zero. This proves
\eqref{eq:constructed-zero-mean} without asserting ergodicity.

Finally, writing $C=E(0)$ gives $D_\gamma(x)=E(x)-C$.
Equation~\eqref{eq:constructed-zero-mean} therefore gives the Ces\`aro
limit in \eqref{eq:cesaro}, with $L_\rho(\gamma)=-C$.
The particle law is supported on the declared admissibility domain.
Canonical replacement leaves the reconstructed exterior field
unchanged, so that exterior particles and their retained interior
count determine the boundary value. Its field-DLR kernel is thus
exactly the particle kernel of Definition~\ref{def:particle-dlr}.
\end{proof}

Combining Proposition~\ref{prop:existence} with
Theorem~\ref{thm:field-unique} and the lifting theorem completes
existence and uniqueness for the particle specification. The moments
and zero-mean property used in this section were obtained for the
constructed law; they were not hypotheses on the competing laws in
the uniqueness theorem.

\section{Consequences and relation to the literature}\label{sec:discussion}

\begin{proof}[Proof of Theorem~\ref{thm:particle-main}]
Corollary~\ref{cor:cesaro-lift} maps every candidate to a jointly
stationary field-DLR law. Theorem~\ref{thm:field-unique} identifies
all such lifts and hence their particle marginals. Proposition~\ref{prop:existence}
supplies a candidate for the declared summation convention. The
almost-sure density statement is part of Theorem~\ref{thm:lift}.
\end{proof}

\begin{corollary}[All finite-intensity stationary Gaussian-kernel laws]
Let $P$ be stationary and satisfy the conditional-density hypothesis of
Theorem~\ref{thm:lift} with finite intensity $\lambda\ge0$.
Then $0\le\lambda\le\rho$ and
\[
P=\frac{\lambda}{\rho}P_{\beta,\rho}
+\left(1-\frac{\lambda}{\rho}\right)\delta_{\varnothing}.
\]
\end{corollary}
\begin{proof}
For $\lambda>0$, Theorem~\ref{thm:lift} determines the empty mass,
and Theorem~\ref{thm:field-unique} identifies the lifted nonempty
component with the constructed law. A stationary zero-intensity
counting measure is empty almost surely.
\end{proof}

This corollary underscores why mean intensity cannot simply be replaced
by an unproved almost-sure density assertion. The extra canonical vacuum
component is a real possibility when intensity is not fixed to the
background. Conversely, no random conditional distribution of the
additive field constant is needed for the particle laws treated here:
their finite Gaussian conditional coefficients already determine it.

For the unique law, the moment estimates imported only for existence
now hold automatically. In particular, the construction and uniqueness
give $\E E(0)^2<\infty$ and
\[
\Var N_{(0,t]}=\E(E(t)-E(0))^2\le4\E E(0)^2.
\]
These are conclusions, not requirements placed on competitors.

Kunz's transfer-operator analysis~\cite{Kunz} and the field viewpoint
of Aizenman--Martin~\cite{AM} are classical foundations of this model.
The published overview~\cite[Section~VI.A.2.e]{Lewin} describes the
periodic crystal states, their uniform phase average, and the
renormalized canonical equations. We do not use that overview's
classification statement as a theorem covering the present candidate
class: the summation convention and the conditional lifting argument
are supplied explicitly above. Nor do we assign an unverified exact
admissibility class to either of the original classical papers.

The reconstruction discussion in \cite[Section~5, equations~(30)--(32)
and the ensuing remarks]{AJJ} uses estimates for finite-volume limits.
It does not by itself establish those estimates for every arbitrary
particle-DLR solution. Theorem~\ref{thm:lift} supplies that missing
logical implication here through nested conditional densities.

Crystalline phase multiplicity is compatible with the theorem. A
fixed-phase law need not be invariant under all real translations; a
uniform average over a periodic phase family is stationary. No
classification of all nonstationary phases was used at any step of the
uniqueness proof. The scope is precisely the finite-coefficient
conditional-density class of Theorem~\ref{thm:lift}, and, for the
existence-and-uniqueness assertion, the particle summation convention
of Definition~\ref{def:particle-dlr}.

\begin{thebibliography}{9}

\bibitem{AM}
M.~Aizenman and P.~A.~Martin,
\emph{Structure of Gibbs states of one dimensional Coulomb systems},
Communications in Mathematical Physics \textbf{78} (1980), 99--116.
\href{https://doi.org/10.1007/BF01941972}{doi:10.1007/BF01941972}.

\bibitem{AJJ}
M.~Aizenman, S.~Jansen and P.~Jung,
\emph{Symmetry breaking in quasi-1D Coulomb systems},
Annales Henri Poincar\'e \textbf{11} (2010), 1453--1485.
\href{https://arxiv.org/abs/1008.3678}{arXiv:1008.3678};
\href{https://doi.org/10.1007/s00023-010-0067-y}{doi:10.1007/s00023-010-0067-y}.
Numbered references in the text use arXiv version~2.

\bibitem{BGL}
D.~Bakry, I.~Gentil and M.~Ledoux,
\emph{Analysis and Geometry of Markov Diffusion Operators},
Grundlehren der mathematischen Wissenschaften, vol.~348,
Springer, Cham, 2014.
\href{https://doi.org/10.1007/978-3-319-00227-9}{doi:10.1007/978-3-319-00227-9}.

\bibitem{Kunz}
H.~Kunz, \emph{The one-dimensional classical electron gas},
Annals of Physics \textbf{85} (1974), 303--335.
\href{https://doi.org/10.1016/0003-4916(74)90413-8}{doi:10.1016/0003-4916(74)90413-8}.

\bibitem{LastPenrose}
G.~Last and M.~Penrose, \emph{Lectures on the Poisson Process},
Institute of Mathematical Statistics Textbooks, vol.~7,
Cambridge University Press, Cambridge, 2017.
\href{https://doi.org/10.1017/9781316104477}{doi:10.1017/9781316104477}.

\bibitem{Lewin}
M.~Lewin, \emph{Coulomb and Riesz gases: The known and the unknown},
Journal of Mathematical Physics \textbf{63} (2022), 061101.
\href{https://arxiv.org/abs/2202.09240}{arXiv:2202.09240};
\href{https://doi.org/10.1063/5.0086835}{doi:10.1063/5.0086835}.

\end{thebibliography}
\end{document}
